\chapter*{General Conclusion and Perspectives}
\phantomsection
\addcontentsline{toc}{chapter}{General Conclusion and Perspectives}
\label{chap:cascade-general-conclusion}

This internship produced one standalone product-comparison assignment followed by an
independent sequence of benchmark and tooling contributions. Keeping those workstreams
separate avoids implying a methodological dependency that did not exist.

For the board-equivalence assignment, the host team supplied the competitor-board list and
asked us to identify the nearest STM32 counterparts. The resulting mappings are useful only
when their criteria are stated, particularly where
the paired boards use different Cortex-M architectures or expose different memory,
peripheral, and software environments. The official ST product page confirms that the
``SM32H735G-DK'' spelling in the internal summary refers to
STM32H735G-DK~\cite{st-stm32h735g-dk}. This assignment did not
select or constrain the boards used in the later benchmark chapters.

The independent benchmark work began with the USB pilot. Optimising the STM32H573I-DK reference
reduced RAM and initialisation time in HID, CDC, and MSC, while ROM remained almost
unchanged in two scenarios. In the recorded Renesas comparison, the optimised STM32 project
has lower summary ROM, RAM, and initialisation values than RA6M5 in all three scenarios.
This is an implementation-level result, not a universal vendor ranking.

The NXP pilot exposed a different trade-off. The workbook lists an average STM32 ROM gap
of $+26.22\,\%$ when RCC is included and $+4.16\,\%$ when it is removed. MCX-N9 retains
the lower summary RAM values and lower recorded initialisation values, while the HID ROM
ordering changes with the measurement boundary. The reviewed presentation resolves the
conflicting workbook cells: its tables confirm the reported HID and CDC ROM values and the
MSC RAM values, while its flowchart identifies the timing case as MSC host initialisation.

The manual reconciliation required by these pilots motivated MCU Workflow Studio and the
MCU Benchmark Copilot Agent. The first structures linker-map analysis and cross-vendor
classification; the second assists project preparation, campaign execution, provenance,
and fairness review. Their value lies in formalising and supporting the method. They do not
replace source artifacts, hardware evidence, or engineering judgement.

The later GD32 and TI campaigns broadened the experimental scope. The GD32 work combined
USB and FreeRTOS observations and revealed different flash, RAM, and reported timing
trade-offs rather than a uniform ecosystem advantage. The TI work combined static driver
analysis, matched bare-metal projects, and FreeRTOS scenarios. Across twelve independently
linked bare-metal images, STM32C5 used 2.8\,\% more aggregate flash and 5.0\,\% less
aggregate RAM. In the three FreeRTOS images, STM32C5 used 27.9--29.6\,\% less flash,
whereas MSPM33 retained a 56--76-byte RAM advantage. Unequal clocks and differing kernel
configurations still prevent a scheduler-efficiency conclusion.

\section*{Perspectives}

Runtime experiments should be repeated with documented equal clocks, software versions,
synchronisation primitives, intervals, warm-up policy, and sample counts, with dispersion
reported alongside central values. The internal evidence register should also be completed
with archived map files, project configurations, and hashes for every published result.

Tool validation forms a second workstream. MCU Workflow Studio should be evaluated on a
larger labelled collection of linker maps, while the MCU Benchmark Copilot Agent should be
assessed through recorded campaigns that distinguish successful automation, manual
intervention, and hardware-dependent steps. Finally, the campaign can be extended to new
middleware and vendors only when each addition preserves explicit measurement boundaries
and fairness limitations. These steps would broaden the evidence without weakening the
central rule established by this work: no comparative conclusion should exceed the
measurements and controls that support it.